\section{The Scene Contract and Renderer}
\label{sec:scene}

The \texttt{Scene} is the single rendering contract. Diagram layout engines speak
only \texttt{Scene}; the renderer understands only \texttt{Scene}. Neither knows
about the other's specifics.

\subsection{Scene and its elements}

A \texttt{Scene} (\texttt{src/contracts/scene.ts}) is a viewBox, an optional
background, a flat list of elements, and optional raw \texttt{defs} (e.g.\ arrow
markers):

\begin{lstlisting}[language=ts]
interface Scene {
  viewBox: { x: number; y: number; width: number; height: number };
  background?: Color;
  elements: SceneElement[];
  defs?: string[];
}
\end{lstlisting}

\texttt{SceneElement} is a five-variant union --- the entire vocabulary the
renderer must support:

\begin{center}
\begin{tabular}{ll}
\toprule
Variant & Carries \\
\midrule
\texttt{rect}   & bounds, fill, stroke, stroke width, optional corner radius \\
\texttt{circle} & centre, radius, fill, stroke \\
\texttt{path}   & SVG path \texttt{d}, stroke, fill, dash, markers, animation \\
\texttt{text}   & content, position, font size/family, fill, anchor, weight \\
\texttt{group}  & children, optional transform / id / opacity \\
\bottomrule
\end{tabular}
\end{center}

That deliberately small surface is why one renderer covers every diagram: a tree,
a slotted page, and a flowchart all reduce to rectangles, circles, paths, and
text.

\subsection{The Pen}

Layout code does not build element objects by hand. A \texttt{Pen}
(\texttt{src/scene/build.ts}), bound once to a theme, produces elements with the
theme's font family defaulted and the \texttt{type} discriminant removed:

\begin{lstlisting}[language=ts]
const p = pen(theme);
p.rect(bounds, fill, stroke, strokeWidth, { rx });
p.circle(center, radius, fill, stroke, strokeWidth);
p.path(d, stroke, strokeWidth, { dash, markerEnd, animated });
p.text(content, x, y, size, fill, { anchor, weight });
\end{lstlisting}

\subsection{The renderer}

\texttt{renderSVG(scene)} (\texttt{src/render/svg.ts}) emits the SVG string
directly --- one element at a time, in document order, with XML-escaped text and
fixed numeric formatting. There is exactly one backend. PNG, when needed, is the
same SVG rasterised by \texttt{resvg}; there is no Skia, PPTX, PDF, or HTML
backend.
